\textbf{Of the properties we put to the split, the lineage label sorts the magnitudes but not the sign.} A Kruskal--Wallis test on that label, read against $1000$ permutations of it, gives $H=13.7$ against a null centred at $5.0$ ($\mathrm{frac}(\mathrm{null}\ge H)<0.001$), and the label accounts for $96\%$ of the between-model variance where a random relabelling accounts for $33\%$. \textbf{That is close to the ceiling these group sizes allow}: the best possible labelling of sixteen values into groups of $3,3,2,6,1,1$ reaches only $H=14.0$. \textbf{The label has six levels, not the five lineages of \S\ref{sec:intro}}: Google ($3$), Meta ($2$), AI2 ($3$), Qwen ($6$), Mistral ($1$) and OLMo-3 ($1$). OLMo-3-7B is its own level because its preference stage is a different construction from the T\"ulu-3 mixture the other three AI2 models share: its card names \texttt{Dolci-Think-DPO-7B}, built with the preference heuristic of \citet{geng2025delta}. \textbf{The sign is a weaker reading, and it depends on where OLMo-3 is put}: with OLMo-3 as its own level, $13$ of the $14$ models that carry a sign fall on their lineage's majority side, against $10.3$ expected under relabelling, $\mathrm{frac}=0.030$ --- above the floor of this test and only just inside the bar --- and with OLMo-3 merged into AI2, $12$ of the $14$ do, against $9.9$ expected, $\mathrm{frac}=0.092$, which does not clear it. Five of the six lineages are sign-pure; the exception is a single Qwen2.5 model. The magnitude reading does not depend on where we put OLMo-3: merged into AI2, which is the labelling our own counterexample used, the test still fires ($H=10.4$, $\mathrm{frac}=0.010$), though that merge makes a second lineage mixed in sign. Lineage and pipeline are confounded in this panel, so this sorts the split without saying which of the two does it.
